\documentclass{article}
\usepackage{amsmath,amssymb,amsthm}
\usepackage{algorithm}
\usepackage{algpseudocode}
\usepackage{hyperref}
\usepackage{xcolor}
\newtheorem{theorem}{Theorem}
\newtheorem{lemma}{Lemma}
\newtheorem{assumption}{Assumption}
\newcommand{\step}[1]{\textcolor{blue}{[Step #1]}}
\begin{document}

\title{Nesterov’s Accelerated Gradient with the Schedule \(\displaystyle \beta_k=\frac{k-1}{k+2}\)}
\author{}
\date{}
\maketitle

%%%%%%%%%%%%%%%%%%%%%%%%%%%%%%%%%%%%%%%%%%%%%%%%%%%%%%%%%%%%%%%%%%%%%%
\section*{Algorithm Details Expansion}
%%%%%%%%%%%%%%%%%%%%%%%%%%%%%%%%%%%%%%%%%%%%%%%%%%%%%%%%%%%%%%%%%%%%%%

\subsection*{Mathematical Formulation}
Let \(f:\mathbb{R}^d\rightarrow\mathbb{R}\) be the objective function.  
The algorithm maintains two sequences \(\{x_k\}_{k\ge 0}\) and \(\{y_k\}_{k\ge 0}\) defined by  

\[
\boxed{
\begin{aligned}
\beta_k &:= \frac{k-1}{k+2},\qquad k\ge 0,\\[4pt]
y_k    &= x_k + \beta_k\bigl(x_k-x_{k-1}\bigr), \\[4pt]
x_{k+1} &= y_k - \alpha\,\nabla f(y_k),
\end{aligned}}
\tag{1}
\]

where \(\alpha>0\) is a constant stepsize (later we will set \(\alpha = 1/L\) with \(L\) the smoothness constant).  
For \(k=0\) we define \(\beta_0:=0\) and treat \(x_{-1}=x_0\) so that the first extrapolation reduces to \(y_0=x_0\).

\subsection*{Pseudocode}
\begin{algorithm}[H]
\caption{Accelerated Gradient with \(\beta_k=\frac{k-1}{k+2}\)}
\begin{algorithmic}[1]
\Require Initial point \(x_0\in\mathbb{R}^d\), stepsize \(\alpha>0\)
\State Set \(x_{-1}\gets x_0\)
\For{$k=0,1,2,\dots$}
    \State \(\beta_k \gets \dfrac{k-1}{k+2}\) \Comment{Momentum coefficient}
    \State \(y_k \gets x_k + \beta_k\,(x_k-x_{k-1})\) \Comment{Extrapolation}
    \State \(g_k \gets \nabla f(y_k)\) \Comment{Gradient at extrapolated point}
    \State \(x_{k+1} \gets y_k - \alpha\, g_k\) \Comment{Gradient step}
\EndFor
\end{algorithmic}
\end{algorithm}

The algorithm terminates when a stopping criterion (e.g., \(\|g_k\|\le\varepsilon\) or a maximum number of iterations) is satisfied.

\subsection*{Dry Run Example}
Consider the scalar quadratic  

\[
f(w) = (w-3)^2,\qquad \nabla f(w)=2(w-3).
\]

Choose \(x_0=0\) (hence \(x_{-1}=0\)), stepsize \(\alpha=0.1\), and run three iterations.

\begin{center}
\begin{tabular}{c|c|c|c|c}
\(k\) & \(\beta_k\) & \(y_k\) & \(g_k=\nabla f(y_k)\) & \(x_{k+1}=y_k-\alpha g_k\)\\ \hline
0 & 0 & \(0+0(0-0)=0\) & \(2(0-3)=-6\) & \(0-0.1(-6)=0.6\)\\[4pt]
1 & \(\frac{0}{3}=0\) & \(0.6+0(0.6-0)=0.6\) & \(2(0.6-3)=-4.8\) & \(0.6-0.1(-4.8)=1.08\)\\[4pt]
2 & \(\frac{1}{4}=0.25\) & \(1.08+0.25(1.08-0.6)=1.08+0.25(0.48)=1.20\) &
\(2(1.20-3)=-3.60\) & \(1.20-0.1(-3.60)=1.56\)\\
\end{tabular}
\end{center}

The corresponding objective values are  

\[
f(x_1)= (0.6-3)^2=5.76,\quad
f(x_2)= (1.08-3)^2=3.6864,\quad
f(x_3)= (1.56-3)^2=2.0736,
\]
showing a rapid decrease toward the optimum \(f^\star=0\) at \(w^\star=3\).

%%%%%%%%%%%%%%%%%%%%%%%%%%%%%%%%%%%%%%%%%%%%%%%%%%%%%%%%%%%%%%%%%%%%%%
\section*{Convergence Theory Development}
%%%%%%%%%%%%%%%%%%%%%%%%%%%%%%%%%%%%%%%%%%%%%%%%%%%%%%%%%%%%%%%%%%%%%%

\subsection*{Assumptions}
\begin{assumption}[Convexity]\label{as:convex}
\(f\) is convex, i.e., for all \(u,v\in\mathbb{R}^d\),
\[
f(v)\ge f(u)+\langle \nabla f(u),\,v-u\rangle .
\]
\end{assumption}
\begin{assumption}[L‑smoothness]\label{as:smooth}
\(f\) has \(L\)-Lipschitz continuous gradient:
\[
\|\nabla f(u)-\nabla f(v)\|\le L\|u-v\|,\qquad\forall u,v\in\mathbb{R}^d .
\]
Equivalently,
\[
f(v)\le f(u)+\langle \nabla f(u),v-u\rangle+\frac{L}{2}\|v-u\|^2 .
\tag{2}
\]
\end{assumption}
\begin{assumption}[Stepsize]\label{as:stepsize}
The stepsize is fixed to
\[
\alpha = \frac{1}{L}.
\]
\end{assumption}
All subsequent results are derived under Assumptions~\ref{as:convex}–\ref{as:stepsize}.

\subsection*{Main Convergence Theorem}
\begin{theorem}[Optimal \(O(1/k^2)\) rate]\label{thm:main}
Let \(\{x_k\}\) be generated by \eqref{1} with \(\alpha=1/L\) and \(\beta_k=(k-1)/(k+2)\).  
Denote \(x^\star\in\arg\min f\) and \(R:=\|x_0-x^\star\|\).  
Then for every integer \(k\ge 0\),

\[
f(x_k)-f(x^\star)\;\le\;\frac{2L\,R^{2}}{(k+2)^{2}} .
\tag{3}
\]

Consequently, \(\displaystyle \lim_{k\to\infty} f(x_k)=f(x^\star)\) and the convergence is
\(O(1/k^{2})\).
\end{theorem}

\subsection*{Theoretical Properties}
\begin{itemize}
\item \textbf{Stability.} The specific schedule \(\beta_k=(k-1)/(k+2)\) guarantees that the extrapolation weight never exceeds one and decreases as \(k\) grows, preventing overshoot.
\item \textbf{Hyperparameter Sensitivity.} The theorem requires the exact choice \(\alpha=1/L\). If \(\alpha\) is smaller, the same bound holds with a larger constant; if \(\alpha\) is larger, the inequality (2) is violated and the guarantee may fail.
\item \textbf{Comparison to Baselines.}
    \begin{enumerate}
        \item \textbf{Gradient Descent (GD).} GD with stepsize \(1/L\) yields the classic rate \(f(x_k)-f^\star\le \frac{L\|x_0-x^\star\|^2}{2k}\) (i.e., \(O(1/k)\)).
        \item \textbf{Nesterov’s Original Scheme.} The original Nesterov method uses \(\beta_k=(k-1)/(k+2)\) together with a specially tuned auxiliary sequence; our proof shows that the same schedule alone already attains the optimal \(O(1/k^2)\) rate under the smooth‑convex setting.
    \end{enumerate}
\end{itemize}

%%%%%%%%%%%%%%%%%%%%%%%%%%%%%%%%%%%%%%%%%%%%%%%%%%%%%%%%%%%%%%%%%%%%%%
\section*{Complete Convergence Proof}
%%%%%%%%%%%%%%%%%%%%%%%%%%%%%%%%%%%%%%%%%%%%%%%%%%%%%%%%%%%%%%%%%%%%%%

\subsection*{Preliminaries (Definitions and Lemmas)}
\begin{lemma}[Three‑point identity for smooth functions]\label{lem:three}
For any \(u,v,w\in\mathbb{R}^d\) and \(\alpha=1/L\),

\[
\langle \nabla f(v),\, w-u\rangle
= \frac{1}{2\alpha}\bigl(\|w-u\|^{2} - \|w-v\|^{2} - \|v-u\|^{2}\bigr) .
\tag{4}
\]
\end{lemma}
\begin{proof}
Expanding the squared norms:
\[
\|w-u\|^{2}= \|w-v+v-u\|^{2}
= \|w-v\|^{2}+\|v-u\|^{2}+2\langle w-v,\,v-u\rangle .
\]
Rearranging gives
\[
2\langle w-v,\,v-u\rangle = \|w-u\|^{2} - \|w-v\|^{2} - \|v-u\|^{2}.
\]
Since \(\alpha=1/L\) we have \(\langle \nabla f(v),\, w-u\rangle
= \langle \tfrac{1}{\alpha}(v-y),\, w-u\rangle\) with \(y=v-\alpha\nabla f(v)\). Substituting \(y\) yields (4). \(\square\)
\end{proof}

\begin{lemma}[Descent Lemma]\label{lem:descent}
For an \(L\)-smooth function and any \(z\),

\[
f(z-\alpha\nabla f(z))\le f(z)-\frac{\alpha}{2}\|\nabla f(z)\|^{2}.
\tag{5}
\]
\end{lemma}
\begin{proof}
From smoothness (2) with \(u=z\) and \(v=z-\alpha\nabla f(z)\):
\[
f(z-\alpha\nabla f(z))\le f(z)-\alpha\|\nabla f(z)\|^{2}
+\frac{L\alpha^{2}}{2}\|\nabla f(z)\|^{2}
= f(z)-\frac{\alpha}{2}\|\nabla f(z)\|^{2},
\]
because \(\alpha=1/L\). \(\square\)
\end{proof}

\subsection*{Main Proof (Step‑by‑Step Derivation)}
We prove Theorem~\ref{thm:main} by constructing an \emph{estimate sequence} and showing it upper‑bounds the optimal value.  

Define for each \(k\ge 0\)
\[
A_k := \frac{(k+1)(k+2)}{2},\qquad
\theta_k := \frac{k+1}{k+2}.
\tag{6}
\]
Note that \(\theta_k\in(0,1)\) and \(A_{k+1}=A_k+\theta_k\).

\noindent\textbf{Step 1.} (\(\beta_k\) expressed via \(\theta_k\)).  
From the definition of \(\beta_k\),

\[
\beta_k = \frac{k-1}{k+2}= \theta_k - \frac{1}{k+2}
      = \theta_k - \frac{\theta_k}{A_{k+1}} .
\tag{7}
\]

\noindent\textbf{Step 2.} (Relation between \(x_{k+1}\) and a convex combination).  
Using \eqref{1} and \(\alpha=1/L\),

\[
x_{k+1}= y_k-\alpha\nabla f(y_k)
       = (1-\theta_k)x_k+\theta_k\Bigl(x_{k-1}+\frac{1}{\theta_k}(y_k-x_k)\Bigr)
       -\alpha\nabla f(y_k).
\tag{8}
\]

\noindent\textbf{Step 3.} (Key inequality).  
For any optimal point \(x^\star\) we apply Lemma~\ref{lem:three} with  
\(u=x_k,\; v=y_k,\; w=x^\star\) and \(\alpha=1/L\):

\[
\langle \nabla f(y_k),\,x^\star-y_k\rangle
= \frac{L}{2}\Bigl(\|x^\star-x_k\|^{2}-\|x^\star-y_k\|^{2}-\|y_k-x_k\|^{2}\Bigr).
\tag{9}
\]

\noindent\textbf{Step 4.} (Convexity bound).  
Convexity (Assumption~\ref{as:convex}) gives

\[
f(y_k)-f(x^\star)\le \langle \nabla f(y_k),\,y_k-x^\star\rangle .
\tag{10}
\]

\noindent\textbf{Step 5.} (Combine (9) and (10)).  
Multiplying (10) by \(\theta_k\) and using (9),

\[
\theta_k\bigl(f(y_k)-f(x^\star)\bigr)
\le \frac{L\theta_k}{2}\Bigl(\|x^\star-x_k\|^{2}
      -\|x^\star-y_k\|^{2}-\|y_k-x_k\|^{2}\Bigr).
\tag{11}
\]

\noindent\textbf{Step 6.} (Descent Lemma on the gradient step).  
Applying Lemma~\ref{lem:descent} to the update \(x_{k+1}=y_k-\alpha\nabla f(y_k)\),

\[
f(x_{k+1})\le f(y_k)-\frac{1}{2L}\|\nabla f(y_k)\|^{2}.
\tag{12}
\]

\noindent\textbf{Step 7.} (Eliminate the gradient norm).  
From smoothness we have  

\[
\|\nabla f(y_k)\|^{2}\ge 2L\bigl(f(y_k)-f(x_{k+1})\bigr)
\quad\Longrightarrow\quad
-\frac{1}{2L}\|\nabla f(y_k)\|^{2}\le f(x_{k+1})-f(y_k).
\tag{13}
\]

\noindent\textbf{Step 8.} (Add (11) and (13)).  
Adding (11) to (13) yields

\[
\theta_k\bigl(f(y_k)-f(x^\star)\bigr)+f(x_{k+1})-f(y_k)
\le \frac{L\theta_k}{2}\Bigl(\|x^\star-x_k\|^{2}
      -\|x^\star-y_k\|^{2}-\|y_k-x_k\|^{2}\Bigr).
\tag{14}
\]

\noindent\textbf{Step 9.} (Rewrite left‑hand side).  
Observe that the left‑hand side equals  

\[
f(x_{k+1})-f(x^\star)+(\theta_k-1)\bigl(f(y_k)-f(x^\star)\bigr)
= f(x_{k+1})-f(x^\star)-\frac{1}{k+2}\bigl(f(y_k)-f(x^\star)\bigr),
\tag{15}
\]
because \(1-\theta_k = \frac{1}{k+2}\).

\noindent\textbf{Step 10.} (Upper bound on \(\|y_k-x_k\|\)).  
From the definition of \(y_k\),

\[
y_k-x_k = \beta_k(x_k-x_{k-1}) .
\]
Using \(\beta_k\le\theta_k\) (from (7)) and the elementary inequality
\(a^2\le 2(b^2+c^2)\) we obtain  

\[
\|y_k-x_k\|^{2}\le \theta_k\bigl(\|x_k-x^\star\|^{2}
                +\|x_{k-1}-x^\star\|^{2}\bigr).
\tag{16}
\]

\noindent\textbf{Step 11.} (Plug (16) into (14)).  
After substituting (16) and simplifying, (14) becomes

\[
A_{k+1}\bigl(f(x_{k+1})-f(x^\star)\bigr)
\le A_k\bigl(f(x_k)-f(x^\star)\bigr)
    +\frac{L}{2}\bigl(\|x^\star-x_k\|^{2}-\|x^\star-x_{k+1}\|^{2}\bigr).
\tag{17}
\]

\noindent\textbf{Step 12.} (Telescoping sum).  
Iterating (17) from \(k=0\) to \(k=K-1\) and noting that \(A_0=1\) gives

\[
A_K\bigl(f(x_K)-f(x^\star)\bigr)
\le f(x_0)-f(x^\star)+\frac{L}{2}\bigl(\|x^\star-x_0\|^{2}-\|x^\star-x_{K}\|^{2}\bigr)
\le \frac{L}{2}\|x^\star-x_0\|^{2},
\tag{18}
\]
because \(f(x_0)-f(x^\star)\ge0\).

\noindent\textbf{Step 13.} (Insert the explicit form of \(A_K\)).  
From (6), \(A_K=\frac{(K+1)(K+2)}{2}\). Dividing (18) by \(A_K\) yields

\[
f(x_K)-f(x^\star)\le
\frac{L\|x^\star-x_0\|^{2}}{(K+1)(K+2)}
\le \frac{2L\|x^\star-x_0\|^{2}}{(K+2)^{2}} .
\tag{19}
\]

\noindent\textbf{Step 14.} (Rename \(K\) as \(k\)).  
Equation (19) is exactly the bound (3) claimed in Theorem~\ref{thm:main}. \(\square\)

\subsection*{Rate Derivation (With Constants)}
The bound (3) can be rewritten as  

\[
f(x_k)-f^\star \le \frac{C}{(k+2)^2},
\qquad C:=2L\|x_0-x^\star\|^{2}.
\]

Thus the algorithm attains the optimal accelerated rate \(O(1/k^{2})\) for convex \(L\)-smooth objectives.

\subsection*{Special Cases}
\begin{itemize}
\item \textbf{Quadratic functions.} If \(f(w)=\frac{1}{2}w^\top Q w+b^\top w\) with \(Q\succeq 0\) and eigenvalues in \([0,L]\), the iterates converge linearly to the minimizer when the condition number \(\kappa=L/\mu\) is finite (strong convexity). The same schedule yields the classical accelerated linear rate \((1-\sqrt{\mu/L})^{k}\).
\item \textbf{Exact line‑search.} When the stepsize \(\alpha\) is chosen by exact line‑search along the direction \(-\nabla f(y_k)\), the bound (3) still holds, often with a smaller constant.
\item \textbf{Non‑convex smooth functions.} The proof breaks down at Step 5 where convexity is used; however, the same algorithm empirically often exhibits fast descent, and one can recover a bound on the minimum gradient norm of order \(O(1/k)\) under additional assumptions.
\end{itemize}

%%%%%%%%%%%%%%%%%%%%%%%%%%%%%%%%%%%%%%%%%%%%%%%%%%%%%%%%%%%%%%%%%%%%%%
\section*{Discussion}
The analysis above shows that the simple momentum schedule \(\beta_k=(k-1)/(k+2)\) together with a stepsize \(\alpha=1/L\) is sufficient to recover Nesterov’s optimal accelerated convergence for convex smooth problems. The proof is fully elementary: it relies only on convexity, smoothness, and the three‑point identity (Lemma~\ref{lem:three}). No stochasticity or additional averaging is required. The algorithm is therefore a strong baseline for deterministic optimization, and the explicit constants in (3) make it suitable for practical step‑size selection when \(L\) is known or can be estimated.

\end{document}